% ============================================================
%  DiffGenMed-VQA — Poster Sunum Konuşma Metni
%  CENG416 Graduation Poster — 2. Kolon (Method & Results)
%  Hedef süre: ~1 dakika
% ============================================================
\documentclass[12pt, a4paper]{article}

\usepackage[utf8]{inputenc}
\usepackage[T1]{fontenc}
\usepackage{lmodern}
\usepackage{geometry}
\usepackage{setspace}

\geometry{margin=2.5cm}
\onehalfspacing

\title{\textbf{DiffGenMed-VQA — Sunum Konuşma Metni}\\
\large 2. Kolon: Method \& Architecture + Results}
\author{Ali Ekin Özçetin \and Ahmet Barış Özel \and Mert Güner}
\date{Haziran 2026}

\begin{document}
\maketitle

% ============================================================
\section*{English (speak this)}
% ============================================================

\begin{quote}
\large

We started with one research question:
\textit{``Can a diffusion-based generative model produce accurate free-form answers to medical visual questions?''}

In our architecture, images are encoded with CLIP, and we compared three language encoders — BERT, BioBERT, and PubMedBERT — to see how pre-training data affects answer quality.
The signals from both encoders are combined in our Cross-Modal Fusion module with learned dynamic weights.
Then, using CIGN, we condition both the forward and reverse diffusion with multimodal features to generate answers from pure noise.

The table shows the key hyperparameters we used in training and sampling.

We trained on Kvasir-VQA, a gastrointestinal endoscopy dataset — the split details are in the table on the poster.

The scores may look low — but that is not the full picture.
Unlike other diffusion models, we do not give any hint to the model.
We apply full Gaussian noise, and the model generates answers completely from pure noise.
This is harder, but more honest.

Looking at the bar chart: all three models score above 0.80 on BERTScore, which measures semantic similarity.
The lower scores on BLEU and ROUGE come from short answers — in a continuous embedding space, even a small rounding error can land on the wrong token.

And PubMedBERT outperforms the others across all metrics.
This confirms our hypothesis: an encoder pre-trained only on biomedical text handles Gaussian noise better.

\end{quote}

% ============================================================
\section*{Türkçe (referans için)}
% ============================================================

\begin{quote}
\large

Tek bir araştırma sorusundan yola çıktık:
\textit{``Difüzyon tabanlı üretici bir model, tıbbi görsel sorulara serbest form cevaplar üretebilir mi?''}

Mimarimizde görüntüler CLIP ile encode edilirken, pre-training verisinin cevap kalitesini nasıl etkilediğini görmek için BERT, BioBERT ve PubMedBERT'i karşılaştırdık.
Her iki encoder'dan gelen sinyaller, Cross-Modal Fusion modülümüzde öğrenilmiş dinamik ağırlıklarla birleştirildi.
Ardından CIGN ile hem ileri hem geri difüzyonu multimodal özelliklerle koşullandırarak saf gürültüden cevap ürettik.

Tabloda eğitim ve örneklemede kullandığımız temel hiperparametreler yer alıyor.

Gastroenteroloji endoskopi veri seti olan Kvasir-VQA üzerinde eğittik — bölünme detayları posterdeki tabloda görülebilir.

Skorlar düşük görünebilir — ama tam resim bu değil.
Diğer difüzyon modellerinin aksine modele hiçbir ipucu vermiyoruz — hiç.
Tam Gaussian gürültü uyguluyoruz ve model cevabı tamamen sıfırdan üretiyor.
Bu daha zor ama daha dürüst bir kurulum.

Bar chart'a bakıldığında: üç model de semantik benzerliği ölçen BERTScore'da 0.80 üzerinde çıkıyor.
BLEU ve ROUGE'daki düşük skorlar ise kısa cevaplardan kaynaklanıyor — sürekli embedding uzayında küçük bir rounding hatası yanlış token'a düşürebilir.

PubMedBERT tüm metriklerde diğerlerinden üstün performans gösterdi.
Bu hipotezimizi doğruluyor: yalnızca biyomedikal metinle pre-train edilmiş bir encoder, Gaussian gürültüyle çok daha iyi başa çıkıyor.

\end{quote}

\end{document}
